\section{Design Decisions and Production Concerns}
\label{sec:levels}

Descriptions of harnesses mix three kinds of statement. ``It has a loop'' names a primitive. ``It gives the model forty tools'' names a decision within one. ``It runs commands in a sandbox'' names a measure taken for a reason that applies to every primitive. Keeping the three apart is the main use of the vocabulary. A primitive is a mechanism every harness has, defined by its function on the request path; a harness has exactly five. The other two levels are where harnesses differ and where they are made fit to run.

\subsection{Design decisions}
A design decision is one way of implementing a primitive. It is the level at which harnesses differ, and the level at which they are usually described. A design decision never adds a primitive; it fills one. The test for whether something is a design decision is whether a different harness could make a different choice and still be a harness. Table~\ref{tab:decisions} lists the decisions found in the deconstruction of \S\ref{sec:deconstruction}, by primitive, with the choices the harnesses made.

\begin{table}[t]
\centering
\small
\setlength{\tabcolsep}{4pt}
\caption{Design decisions by primitive, and the choices seen across the harnesses deconstructed in \S\ref{sec:deconstruction}.}
\label{tab:decisions}
\begin{tabularx}{\linewidth}{@{}l>{\raggedright\arraybackslash}p{3.1cm}Y@{}}
\toprule
Primitive & Decision & Choices seen \\
\midrule
\prim{Input} & source & terminal; editor; chat gateways; HTTP; issue tracker; scheduled trigger; a typed state channel \\
 & mid-turn messages & delivered at the next step boundary (steer); after the turn (queue); as an interrupt; not at all \\
 & the person's answers & consumed by the harness as approvals; returned to the model as tool results \\
\addlinespace
\prim{Context} & instructions & a system prompt alone; plus project files (\texttt{CLAUDE.md}, \texttt{AGENTS.md}); plus memory files; a sectioned prompt diffed into patches \\
 & long-term memory & none; files preloaded into the instructions; a store searched by the agent through tools; summaries carried forward \\
 & window fitting & drop old turns; summarize (prune, then summarize); compaction at a threshold with a verbatim tail; provider-native compaction; left to user code \\
 & layout for caching & a stable prefix with changes appended; instruction changes frozen as dated events; byte-stable tiers \\
\addlinespace
\prim{Model interface} & client & provider SDK; in-house request compiled per protocol; one stream function over adapters; a library's chat-model call \\
 & transport & HTTP with server-sent events; WebSocket, falling back to HTTP \\
 & retries and backup & none; three to ten retries; retry split by whether output started; a fallback chain of models \\
 & action format & native tool calls; calls parsed from text; a code block to execute \\
\addlinespace
\prim{Output} & tool set & one shell; a fixed built-in set; built-in plus MCP; plain tool objects; a tool node \\
 & action format & structured calls; model-written code run in an interpreter (V8, QuickJS, in-process) \\
 & execution environment & the host; an OS sandbox with a domain allowlist; a container; a remote sandbox \\
 & parallelism & sequential; read-only calls concurrent; all calls concurrent as they stream \\
\addlinespace
\prim{Control flow} & loop shape & until no tool call; one task per turn; a user-written graph run as supersteps; a phase-verdict loop \\
 & stopping rules & step budget; token or cost budget; unlimited; recursion limit; stop-gate nudges \\
 & approval & a rule pipeline ending in a prompt; a classifier model; a reviewer model; each tool asserting its own permission \\
 & delegation & an agent tool to a fixed depth; child sessions; sub-agents as full sessions; flat threaded delegation \\
\bottomrule
\end{tabularx}
\end{table}

\subsection{Production concerns}
A production concern is a question asked of every primitive, whose answer changes how some primitives are implemented. We use six (Table~\ref{tab:concerns}): what should outlast the session (persistence); what could do harm, and who could make it (safety and security); what did it do (observability); what if something breaks (resilience); what is slow or costly (performance); and how do we know it still works (evaluation). A concern is not a part of the harness. A sandbox is not a sixth primitive; it is the safety answer in \prim{Output}. A memory store is not a sixth primitive; it is a store in the world that \prim{Context} reads and \prim{Output} writes, kept for the persistence concern. A trace is not a sixth primitive; it is the observability answer, recorded wherever the primitive it watches runs. The test for whether something is a concern is whether it is motivated by a question that could equally be asked of a different primitive.

\begin{table}[t]
\centering
\scriptsize
\setlength{\tabcolsep}{3pt}
\caption{Six production concerns as questions asked of each primitive. Each cell names what the concern changes in that primitive; ``---'' means we found nothing essential there in the harnesses examined.}
\label{tab:concerns}
\begin{tabularx}{\linewidth}{@{}l*{6}{Y}@{}}
\toprule
& \textbf{Persistence} & \textbf{Safety \& security} & \textbf{Observability} & \textbf{Resilience} & \textbf{Performance} & \textbf{Evaluation} \\
& \emph{What outlasts the session?} & \emph{What could do harm, and who could make it?} & \emph{What did it do?} & \emph{What if something breaks?} & \emph{What is slow or costly?} & \emph{Does it still work?} \\
\midrule
\prim{Input} & an unfinished session offered back & approvals asked of a person & refusals recorded & input ending; interrupts & --- & yes to every approval in a test run \\
\prim{Context} & memories; resumed transcripts & untrusted content; poisoned memory; roles and authority & --- & compaction when too long & caching the unchanged prefix; trimming results & the instructions under test \\
\prim{Model interface} & --- & --- & tokens, latency, cost per call & retries; a backup model & streaming & the model under test \\
\prim{Output} & --- & sandboxes; never-lists; secrets hidden; pinned images & what each action did & time limits on actions & parallel actions; smaller models for side tasks & the tools under test \\
\prim{Control flow} & the task in progress; a diary of every message & runaway loops; the record kept out of reach & each decision & resume after a crash & --- & the loop under test \\
\bottomrule
\end{tabularx}
\end{table}

The concerns interact. Persistence creates a security surface: a poisoned memory steers every later session. Streaming changes safety: an interrupt can stop a reply mid-sentence, and the harness must record what was said rather than a finished answer. Summarizing changes persistence: the summary must say where the originals are kept. Treating concerns as questions asked of each primitive makes these interactions visible; treating them as separate components hides them.

\subsection{Bundles and the cases at the boundaries}
Much of what a harness's documentation names is a bundle. A mode, a skill, a hook system or an agent profile packages decisions across several primitives: a planning mode is a context template, a restricted tool set and an approval gate; a hook system is a dispatcher in \prim{Control flow} whose events each affect one primitive. Decomposed, each element has one home. The negative definitions decide most placements. Thirteen rules, each needed for a component found in \S\ref{sec:deconstruction}, decide the rest, and are given in Appendix~\ref{app:rules}. Three are used often enough to state here. A component is assigned by its function, not its location: an approval decision made inside a tool's runtime is \prim{Control flow}. A sub-agent is a harness of its own, and delegation to it is \prim{Control flow} of the outer harness. And one primitive may use another's mechanism to do its job: \prim{Context} may ask the model to summarize old turns before compacting them, and the component is assigned by what it decides, with the call it makes a use of \prim{Model interface}.
